\section{Análisis de casos: rutas de ejecución de extremo a extremo para 5 tipos de tareas típicas}

\subsection{Caso 1: análisis de formatos de datos desconocidos y generación de informes de análisis}

Este es el prototipo más representativo de las «tareas de cola larga» en la entrevista. La entrada suele ser un formato de archivo no estándar, documentación dispersa y restricciones incompletas. Una cadena de ejecución estable de un Agent debería ser: identificar el formato -> buscar la herramienta de análisis -> validar la corrección del análisis -> generar datos estructurados intermedios -> producir un informe legible. A diferencia de los scripts tradicionales, la ventaja de un Agent general está en poder corregir la estrategia mientras ejecuta.

\begin{importantbox}{Claves del éxito de este tipo de tarea}
\begin{itemize}
    \item No considerar exitoso el primer análisis como el final; es necesario hacer una validación de consistencia
    \item Distinguir con claridad entre «falla de análisis» y «falla de comprensión semántica»
    \item La salida del informe debe incluir resultados intermedios rastreables, para reducir el costo de la revisión
\end{itemize}
\end{importantbox}

Si se coloca esta tarea en el contexto de la entrevista, puede entenderse como una muestra de la capacidad central de Manus: frente a una distribución de entradas nunca antes vista, el sistema no depende de plantillas fijas, sino que forma soluciones temporales mediante Tool Use y la ejecución de código.

\subsection{Caso 2: integración de información de múltiples fuentes y memorandos de decisión}

Otro tipo de tarea frecuente es «converger la información de múltiples fuentes en un texto que permita tomar una decisión». Lo que esta tarea exige del Agent no es «escribir extenso», sino «información coherente, estructura clara y conclusiones rastreables». Los problemas comunes en la ejecución real incluyen: conflictos entre fuentes sin marcar, mezcla de hechos y opiniones, y ruptura de la cadena de citas.

\begin{warningbox}{Los tres tipos de errores más frecuentes en las tareas de integración de información}
\begin{enumerate}
    \item Tomar una opinión de segunda mano como un hecho de primera mano
    \item La cita existe, pero no puede ubicarse en el contexto original
    \item La fuerza de la conclusión no corresponde a la fuerza de la evidencia
\end{enumerate}
\end{warningbox}

La entrevista enfatiza que «lo que el usuario finalmente quiere es una decisión que le ahorre esfuerzo», por lo que la evaluación de este tipo de tarea no puede fijarse solo en la fluidez del texto, sino también en si la conclusión puede volver a revisarse. En la plantilla de salida se pueden incorporar de manera obligatoria los campos de «nivel de evidencia» y «aviso de contraejemplos».

\subsection{Caso 3: tareas de investigación (Deep Research) y entregas reproducibles}

Las tareas de investigación suelen malinterpretarse como «buscar un poco más de información». En realidad, se acercan más a un mini research pipeline: descomposición del problema, búsqueda de información, filtrado de evidencia, corrección de hipótesis y entrega de conclusiones. El juicio de la entrevista sobre Deep Research es claro: se homogeneizará de manera gradual, por lo que el punto de diferenciación está en la estabilidad de la ejecución y la calidad de la entrega.

\begin{knowledgebox}{Recomendación de estructura de salida para tareas de investigación}
\begin{itemize}
    \item Resumen de conclusiones: de 3 a 5 juicios accionables
    \item Panel de evidencia: cada juicio con su fuente, fecha y nivel de confiabilidad correspondientes
    \item Declaración de incertidumbre: qué conclusiones requieren verificación adicional
    \item Siguiente acción: delimitar con claridad la división del trabajo entre humano y máquina
\end{itemize}
\end{knowledgebox}

Si esta estructura se fija en el producto, puede reducirse de manera significativa el riesgo de una salida que «parece completa, pero en realidad no permite decidir». La filosofía de ingeniería de la entrevista es precisamente esa: convertir una capacidad ambigua en una entrega estructurada.

\subsection{Caso 4: tareas de orquestación de cadenas de herramientas (mezcla de código, navegador y línea de comandos)}

Este tipo de tarea es donde mejor se manifiesta el valor de un «espacio de acción unificado». Los flujos tradicionales suelen separar la automatización del navegador, la ejecución de scripts y la consolidación de resultados en varios sistemas, lo que provoca una ruptura del estado. La ventaja de la ruta de Manus está en converger de manera unificada las acciones dentro del sandbox de la sesión, reduciendo la pérdida de contexto.

\begin{importantbox}{Plantilla de orden de ejecución para tareas de cadenas de herramientas mixtas}
\begin{enumerate}
    \item Definir el estado objetivo (qué se necesita entregar al final)
    \item Planear la ruta viable más corta (el menor número de herramientas y de pasos)
    \item Validar de inmediato después de ejecutar (en lugar de aceptar todo al final de una sola vez)
    \item Revertir de manera parcial en caso de falla (conservando los resultados intermedios ya validados)
\end{enumerate}
\end{importantbox}

El énfasis de la entrevista en «eliminar herramientas» resulta especialmente importante aquí. Cuantas más herramientas existan, mayor es la probabilidad de errores de enrutamiento y de desviación del estado. En vez de ampliar la superficie de capacidades, es preferible mejorar la confiabilidad de la ruta predeterminada.

\subsection{Caso 5: ejecución en segundo plano de tareas largas y estrategia de intervención humana}

Las tareas largas son el escenario de mayor valor y, a la vez, más frágil para un Agent general. La «ejecución continua del sandbox» de la que habla la entrevista resuelve el problema de la continuidad de la sesión, pero además es necesario delimitar con claridad los límites de la intervención humana. De lo contrario, ante una anomalía, la tarea reintentará de manera indefinida o fallará en silencio.

\begin{warningbox}{Las tres reglas de intervención que las tareas largas deben definir de manera explícita}
\begin{itemize}
    \item Cuándo reintentar de manera automática y cuántas veces reintentar
    \item Cuándo debe solicitarse confirmación humana
    \item Cuándo debe activarse una interrupción de seguridad que conserve el estado
\end{itemize}
\end{warningbox}

En cuanto faltan este tipo de reglas, el sistema presenta una automatización aparente que «parece automatizada, pero en realidad es incontrolable». La observabilidad y el mecanismo de reproducción de la entrevista existen precisamente para resolver este problema.

\subsection{Comparación entre casos: capacidades, riesgos y acciones de gobernanza}

\begin{center}
\begin{tabular}{p{3.0cm}p{3.5cm}p{4.1cm}p{3.5cm}}
\toprule
Tipo de tarea & Capacidad clave & Riesgo principal & Acción de gobernanza \\
\midrule
Análisis de formato desconocido & Tool discovery + coding & Análisis correcto pero con lectura semántica errónea & Validación de doble capa (estructura + semántica) \\
Memorando de integración de información & Consolidación de evidencia + redacción estructurada & Ruptura de la cadena de citas & Campos de evidencia obligatorios \\
Tarea de investigación & Descomposición del problema + filtrado de evidencia & Extrapolación excesiva de la fuerza de la conclusión & Declaración de incertidumbre en la salida \\
Cadena de herramientas mixta & Gestión de estado + capacidad de reversión & Enrutamiento erróneo de herramientas & Lista blanca de herramientas y ruta predeterminada \\
Ejecución de tareas largas & Ejecución continua + monitoreo y alertas & Falla silenciosa o reintento indefinido & Reglas de intervención y condiciones de interrupción \\
\bottomrule
\end{tabular}
\end{center}

\subsection{Resumen de la sección}

Los 5 tipos de casos apuntan en conjunto a la misma conclusión: el valor de un Agent general no está en «si sabe hacerlo o no», sino en «si puede terminarlo de manera estable y entregarlo de forma rastreable». Esta es también la razón fundamental por la que la entrevista insiste una y otra vez en la ingeniería de sistemas y la disciplina organizacional.

